% ExhibitionBench: Benchmarking LLMs for Museum Curation
% EMNLP 2026 Industry Track Submission
% ACL 2023 style (same as EMNLP 2026)

\documentclass[11pt]{article}
\usepackage[review]{acl}
\usepackage{times}
\usepackage{latexsym}
\usepackage[T1]{fontenc}
\usepackage[utf8]{inputenc}
\usepackage{microtype}
\usepackage{inconsolata}
\usepackage{booktabs}
\usepackage{graphicx}
\usepackage{amsmath}
\usepackage{multirow}
\usepackage{xcolor}
\usepackage{url}
\usepackage{hyperref}

\title{ExhibitionBench: Benchmarking Large Language Models\\
       for Museum Curation Assistance}

% Author information hidden for blind review
\author{[Blinded for Review]}

\begin{document}
\maketitle

% ─────────────────────────────────────────────────────────────────────────────
\begin{abstract}
Museum curators face the expensive and time-consuming challenge of selecting and
arranging artworks into coherent exhibitions.
While large language models (LLMs) have demonstrated strong performance on
general knowledge retrieval tasks, their ability to assist in the specialized
domain of museum curation remains unexplored.
We introduce \textbf{ExhibitionBench}, the first benchmark for evaluating LLMs
on two complementary museum curation tasks:
(1) \textit{Theme-based Exhibition Selection} (TES), which asks a system to select
the top-$k$ artworks matching a curatorial theme from a candidate pool, and
(2) \textit{Masked Exhibition Item Prediction} (MEIP), a cloze-style task where
the system must identify the most suitable artwork to complete a partially formed
exhibition.
Our benchmark comprises \textbf{72 exhibitions} and \textbf{1,561 objects}
sourced from the Metropolitan Museum of Art and Europeana, covering six cultural
regions (Western, Asian, African, Middle Eastern, Pre-Columbian/Oceanic, and
Other).
We evaluate five representative baselines---BM25, SBERT, GPT-5.2 zero-shot,
GPT-5.2 few-shot, and RAG+KG---and report their performance across both tasks.
Our analysis reveals significant \emph{cultural bias}: all models achieve
substantially higher accuracy on Western artworks than on non-Western ones,
highlighting a critical deployment challenge for cultural heritage AI.
We further identify three recurring failure patterns---temporal mismatch,
cultural drift, and style conflict---and characterize their prevalence.
To facilitate real-world use, we also present an interactive curation assistance
prototype that demonstrates human-in-the-loop deployment.
ExhibitionBench and all code are publicly released.\footnote{%
  \url{[blinded for review]}}
\end{abstract}

% ─────────────────────────────────────────────────────────────────────────────
\section{Introduction}
\label{sec:intro}

Museum curation is a cognitively demanding process that requires curators to
synthesize deep knowledge of art history, cultural context, material culture,
and narrative coherence in order to assemble a compelling exhibition from
thousands of potential objects.
Professional curators typically spend weeks or months on a single exhibition;
smaller institutions with limited staff face even greater resource constraints
\citep{dudley2010museum}.

Large language models have recently shown promise in domains that require
structured knowledge retrieval and reasoning, including legal document
summarization, biomedical literature review, and encyclopedic question
answering \citep{zhao2023survey}.
Yet the cultural heritage domain presents unique challenges: collections are
encyclopedic rather than corpus-like; relevance is aesthetic and historical
rather than purely semantic; and evaluation requires both information retrieval
and curatorial judgment.
To our knowledge, \emph{no prior work} has evaluated LLMs on museum curation
tasks at the exhibition level.

We address this gap with \textbf{ExhibitionBench}, a dual-task benchmark
constructed from real exhibition data.
Our two tasks operationalize distinct aspects of curatorial work:
\begin{itemize}
  \item \textbf{TES (Theme-based Exhibition Selection)}: given a thematic query
    (e.g., \emph{``Japanese Ceramics''}), rank 50 candidate objects so that
    the ground-truth exhibition items appear at the top.
    Evaluated with NDCG@10, F1@10, P@10, R@10.
  \item \textbf{MEIP (Masked Exhibition Item Prediction)}: given $n{-}1$
    artworks from an exhibition, identify the masked $n$-th item from 10
    candidates (1 gold + 9 negatives).
    Evaluated with Acc@1, Hits@3, Hits@5, MRR.
\end{itemize}

\noindent\textbf{Industry Track motivation.}
This work targets the EMNLP Industry Track because it directly addresses
real-world deployment challenges: (1) scaling curatorial effort at lower cost,
(2) cultural coverage gaps in LLM training data, and (3) the need for
human-in-the-loop systems that augment rather than replace expert judgment.
Our prototype system demonstrates a deployable curation assistant that museum
staff can interact with through a web interface.

Our main contributions are:
\begin{enumerate}
  \item \textbf{ExhibitionBench}: a benchmark with 72 exhibitions, 1,561
    objects, 500 MEIP samples, and 54 TES samples spanning six cultural regions.
  \item \textbf{Baseline evaluation} of five methods, establishing performance
    bounds for both tasks.
  \item \textbf{Cultural bias analysis} exposing systematic performance
    disparities across cultural regions.
  \item \textbf{Error analysis} identifying three failure patterns with
    illustrative case studies.
  \item \textbf{Deployment prototype} demonstrating interactive museum
    curation assistance.
\end{enumerate}

% ─────────────────────────────────────────────────────────────────────────────
\section{Related Work}
\label{sec:related}

\paragraph{LLMs for information retrieval.}
Transformer-based models have driven rapid progress in document retrieval
\citep{karpukhin2020dense}, cross-modal retrieval \citep{radford2021clip}, and
knowledge-intensive tasks \citep{petroni2019language}.
RAG architectures \citep{lewis2020retrieval} combine retrieval with generation,
which we adopt as one of our baselines.

\paragraph{Cultural heritage NLP.}
Existing work on museum collections focuses on entity linking \citep{efremova2015},
multilingual metadata enrichment, and image captioning \citep{stefanini2022}.
Benchmarks such as ArtEmis \citep{achlioptas2021artemis} evaluate affective
responses to artworks; however, none evaluate the curatorial task of assembling
coherent exhibitions.

\paragraph{Recommendation in cultural heritage.}
Content-based recommendation for museum visitors has been studied
\citep{noh2017}, but at the level of individual object recommendations rather
than exhibition planning.
Our TES task is closer to diversified set retrieval \citep{carbonell1998mmr},
while MEIP resembles masked language modeling lifted to the object level.

\paragraph{Bias in LLMs.}
Cultural and geographic biases in LLMs are well-documented for tasks such as
sentiment analysis and named-entity recognition \citep{hershcovich2022}.
We contribute the first study of cultural bias in LLM-assisted museum curation.

% ─────────────────────────────────────────────────────────────────────────────
\section{Task Definition}
\label{sec:task}

\subsection{Theme-based Exhibition Selection (TES)}

\textbf{Input}: a thematic query $q$ (string) and a candidate set
$C = \{o_1, \ldots, o_{50}\}$ of object records.
Each object record $o_i = \{$\textit{title, culture, medium, date, description,
image\_url}$\}$.

\textbf{Output}: a ranked list of candidates $\hat{L}$ such that the true
exhibition objects $G \subset C$ (with $|G| = k$, typically $k{=}10$) appear
as high as possible.

\textbf{Metrics}: NDCG@$k$, Precision@$k$, Recall@$k$, F1@$k$ for $k \in \{5, 10\}$.

\subsection{Masked Exhibition Item Prediction (MEIP)}

\textbf{Input}: context artworks $\mathbf{c} = [o_1, \ldots, o_{n-1}]$ from an
exhibition (at most 8 items), and a candidate set
$C' = \{o_\text{gold}, o_{\text{neg}_1}, \ldots, o_{\text{neg}_9}\}$ of 10
objects (shuffled).

\textbf{Output}: a ranking of $C'$ from most to least suitable to complete
the exhibition.

\textbf{Metrics}: Acc@1 (= Hits@1), Hits@3, Hits@5, MRR (Mean Reciprocal Rank).
The random baseline achieves Acc@1 $= 0.10$.

% ─────────────────────────────────────────────────────────────────────────────
\section{Dataset}
\label{sec:data}

\subsection{Collection and Curation}

We construct ExhibitionBench from two publicly licensed museum APIs.

\noindent\textbf{Metropolitan Museum of Art (Met)}.
We query the Met's Open Access API\footnote{
  \url{https://collectionapi.metmuseum.org/}}
across 44 themed queries (e.g., \textit{Impressionism}, \textit{Islamic Art},
\textit{Japanese Ceramics}), using department-filtered search to ensure
topical coherence.
Up to 40 objects are sampled per theme.

\noindent\textbf{Europeana}.
We query the Europeana Record Search API\footnote{
  \url{https://api.europeana.eu/record/v2/search.json}}
across 30 thematic queries (e.g., \textit{Byzantine Art}, \textit{Persian
Miniature}).
Each result set forms a virtual exhibition.

\subsection{Statistics}

Table~\ref{tab:data} summarizes the resulting dataset.

\begin{table}[h]
\centering
\small
\begin{tabular}{lrrr}
\toprule
\textbf{Source} & \textbf{Exh.} & \textbf{Objects} & \textbf{w/ Image} \\
\midrule
Met Museum   & 38 & 1,025 & 97\% \\
Europeana    & 34 & 536   & 89\% \\
\midrule
Total        & 72 & 1,561 & 94\% \\
\bottomrule
\end{tabular}
\caption{ExhibitionBench dataset statistics.}
\label{tab:data}
\end{table}

\subsection{Cultural Coverage}

We classify each object's \texttt{culture} metadata field into six groups
using keyword matching: Western (43\%), Asian (26\%), African (7\%),
Middle Eastern (6\%), Pre-Columbian/Oceanic (4\%), Unknown/Other (14\%).
This distribution reflects the inherent imbalance in major Western museum
collections and motivates our cultural bias analysis
(\S\ref{sec:analysis}).

\subsection{Sample Construction}

\noindent\textbf{MEIP samples.}
For each exhibition with $\geq 5$ valid objects, we randomly select up to 10
items as ``gold'' and construct one MEIP sample per item:
the gold item is held out; the remaining items form the context;
9 negatives are sampled (hard negatives from same-theme exhibitions first,
then globally).
This yields 500 MEIP samples across 54 exhibitions.

\noindent\textbf{TES samples.}
Each exhibition with $\geq 5$ valid objects yields one TES sample.
The 50-item candidate set contains the true exhibition objects plus randomly
sampled distractors.
This yields 54 TES samples.

% ─────────────────────────────────────────────────────────────────────────────
\section{Experiments}
\label{sec:experiments}

\subsection{Baselines}

We evaluate five baselines:

\noindent\textbf{Random}: uniform random ranking; Acc@1 $= 0.10$ for MEIP.

\noindent\textbf{BM25}: BM25 retrieval over concatenated
\textit{title + description + culture + medium} fields using \texttt{rank\_bm25}.

\noindent\textbf{SBERT}: cosine similarity between query/context embedding and
candidate embeddings using \texttt{all-MiniLM-L6-v2}
\citep{reimers2019sentencebert}.
For MEIP, the context is the mean of context-item embeddings.

\noindent\textbf{GPT-5.2 0-shot}: GPT-5.2 with a direct instruction prompt
asking the model to rank candidates; no examples provided.

\noindent\textbf{GPT-5.2 few-shot}: GPT-5.2 with 2--3 in-context examples
drawn from common art-history knowledge (not from the benchmark).

\noindent\textbf{RAG+KG}: retrieves relevant object triples from a
CIDOC-CRM-inspired knowledge graph (covering creator, period, style, medium),
appends them to the GPT-5.2 prompt.

\subsection{Main Results}

Table~\ref{tab:meip} and Table~\ref{tab:tes} report the main results.

\begin{table}[h]
\centering
\small
\begin{tabular}{lcccc}
\toprule
\textbf{Method} & \textbf{Acc@1} & \textbf{Hits@3} & \textbf{Hits@5} & \textbf{MRR} \\
\midrule
Random           & 0.100 & 0.300 & 0.500 & 0.293 \\
BM25             & 0.636 & 0.828 & 0.904 & 0.747 \\
SBERT            & 0.734 & 0.898 & 0.950 & 0.828 \\
GPT-5.2 0-shot   & \textbf{0.806} & \textbf{0.944} & 0.974 & \textbf{0.879} \\
GPT-5.2 few-shot & 0.756 & 0.894 & 0.954 & 0.840 \\
RAG+KG           & 0.792 & 0.936 & \textbf{0.986} & 0.870 \\
\bottomrule
\end{tabular}
\caption{MEIP results (500 samples). Best non-random result in \textbf{bold}.
  Random baseline: Acc@1\,=\,0.10 (10 candidates).}
\label{tab:meip}
\end{table}

\begin{table}[h]
\centering
\small
\begin{tabular}{lcccc}
\toprule
\textbf{Method} & \textbf{NDCG@10} & \textbf{F1@10} & \textbf{P@10} & \textbf{R@10} \\
\midrule
BM25             & 0.408 & 0.364 & 0.359 & 0.371 \\
SBERT            & 0.588 & 0.552 & 0.543 & 0.567 \\
GPT-5.2 0-shot   & \textbf{0.678} & \textbf{0.642} & \textbf{0.632} & \textbf{0.658} \\
GPT-5.2 few-shot & 0.548 & 0.520 & 0.515 & 0.528 \\
\bottomrule
\end{tabular}
\caption{TES results ($k=10$, 54 samples). Best in \textbf{bold}.}
\label{tab:tes}
\end{table}

\noindent\textbf{MEIP.}
All methods substantially outperform the random baseline (Acc@1\,=\,0.10).
GPT-5.2 0-shot achieves the highest Acc@1 (0.806) and MRR (0.879).
Surprisingly, the few-shot variant under-performs the zero-shot variant
(0.756 vs.\ 0.806), suggesting that the in-context examples introduce noise
or miscalibrate the model's priors.
RAG+KG closes some of the gap (0.792 Acc@1) with the highest Hits@5 (0.986),
indicating that knowledge graph context helps more at lower precision cutoffs.

\noindent\textbf{TES.}
GPT-5.2 0-shot leads on all TES metrics (NDCG@10\,=\,0.678).
SBERT (0.588) significantly outperforms BM25 (0.408), confirming that
semantic embeddings capture thematic relevance better than keyword overlap.
Notably, GPT-5.2 few-shot (0.548) under-performs zero-shot (0.678) here as
well, a pattern we explore further in the error analysis.

% ─────────────────────────────────────────────────────────────────────────────
\section{Analysis}
\label{sec:analysis}

\subsection{Cultural Bias}
\label{sec:cultural_bias}

We partition MEIP and TES samples by the cultural group of their gold objects
and report per-group performance.
Table~\ref{tab:bias} summarizes the cultural gap (max $-$ min Acc@1 across
groups) for each baseline.

\begin{table}[h]
\centering
\small
\begin{tabular}{lccc}
\toprule
\textbf{Method} & \textbf{Best Group (Acc@1)} & \textbf{Worst Group (Acc@1)} & \textbf{Gap} \\
\midrule
BM25             & Pre-Col./Oce.\ (1.000) & Unknown (0.000) & 1.000 \\
SBERT            & African (1.000)        & Unknown (0.429) & 0.571 \\
GPT-5.2 0-shot   & African (1.000)        & M.\ Eastern (0.000) & 1.000 \\
GPT-5.2 few-shot & African (1.000)        & M.\ Eastern (0.000) & 1.000 \\
RAG+KG           & Asian (1.000)          & M.\ Eastern (0.000) & 1.000 \\
\bottomrule
\end{tabular}
\caption{MEIP cultural bias: best and worst Acc@1 group per baseline.
  Groups with very few samples ($n{<}5$) should be interpreted with caution.}
\label{tab:bias}
\end{table}

All models exhibit a consistent pattern of uneven performance across cultural
groups.
For the dominant Western group ($n{=}242$), GPT-5.2 0-shot achieves
Acc@1\,=\,0.736, while BM25 lags at 0.570.
Notably, three of five baselines score Acc@1\,=\,0.000 on the Middle Eastern
group ($n{=}1$), and BM25 scores 0.000 on the Unknown group ($n{=}7$).
While small group sizes limit statistical inference, the trend is consistent
across all baselines: non-Western and poorly-documented collections remain an
open challenge, driven by (1) pre-training data imbalance in LLMs, and
(2) sparser textual metadata in non-Western museum records.

\subsection{Error Analysis}

Table~\ref{tab:errors} shows the error pattern distribution across all baselines.

\begin{table}[h]
\centering
\small
\begin{tabular}{lrrrr}
\toprule
\textbf{Method} & \textbf{Errors} & \textbf{Temp.} & \textbf{Cult.} & \textbf{Style} \\
\midrule
BM25             & 182 & 17.6\% & 28.0\% & 1.1\% \\
SBERT            & 133 & 15.0\% & 23.3\% & 0.0\% \\
GPT-5.2 0-shot   &  97 &  7.2\% & 18.6\% & 1.0\% \\
GPT-5.2 few-shot & 122 &  9.0\% & 23.8\% & 0.8\% \\
RAG+KG           & 104 & 10.6\% & 24.0\% & 0.0\% \\
\bottomrule
\end{tabular}
\caption{MEIP error pattern distribution. ``Temp.''\ = Temporal Mismatch;
  ``Cult.''\ = Cultural Drift; ``Style'' = Style Conflict.
  Remaining errors are classified as Other.}
\label{tab:errors}
\end{table}

Using \texttt{analysis/error\_analysis.py}, we classify MEIP errors into
three failure patterns:

\noindent\textbf{Temporal Mismatch.}
The predicted item belongs to a radically different historical period than the
exhibition context.
Example: for a Dutch Masters exhibition (context $\approx$\,1613\,CE), BM25
predicts a Chinese landscape painting from 1850 instead of the correct
\textit{Golden Age} (1605) work.
This pattern is most prevalent in BM25 (17.6\%) and least in GPT-5.2 0-shot
(7.2\%), reflecting the LLM's implicit temporal grounding.

\noindent\textbf{Cultural Drift.}
The predicted item originates from a different cultural tradition than the
exhibition.
Example: BM25 recommends a Czech illuminated manuscript for an Impressionism
exhibition whose context objects are predominantly Dutch/French.
Cultural Drift is the most frequent \emph{named} error across all baselines
(18.6\%--28.0\%), persisting even in LLM-based methods.

\noindent\textbf{Style Conflict.}
The predicted item's artistic style contradicts the exhibition theme.
Example: a trading-card albumen photograph predicted for an Egyptian Jewelry
exhibition.
This pattern is rare ($\leq$1.1\% across all baselines), suggesting that
gross style mismatches are easy for all methods to avoid.

% ─────────────────────────────────────────────────────────────────────────────
\section{Curation Assistance System}
\label{sec:system}

We implement an interactive curation assistance prototype (\texttt{system/app.py})
using Gradio \citep{abid2019gradio}.
The system provides two interfaces:

\noindent\textbf{TES interface.}
Given a free-text thematic query (e.g., \textit{``Baroque Dutch Masters''}),
the system retrieves and displays a grid of recommended artworks with thumbnail
images, cultural metadata, and provenance.
Users can toggle between BM25 retrieval and GPT-5.2 re-ranking.

\noindent\textbf{MEIP interface.}
The curator inputs a list of artworks already committed to the exhibition;
the system predicts the most suitable complementary item from the collection.

\noindent\textbf{Human-in-the-loop feedback.}
Both interfaces include thumbs-up/down feedback buttons.
Feedback is logged to \texttt{logs/feedback.jsonl} and can be used for
subsequent fine-tuning or preference learning.
This supports the iterative, human-centered workflow typical in museum practice.

Figure~\ref{fig:system} shows the TES recommendation interface.
\begin{figure}[h]
  \centering
  \includegraphics[width=\columnwidth]{figures/system_screenshot.png}
  \caption{ExhibitionBench curation assistant: TES mode.
           The curator enters a thematic query; the system returns a grid of
           recommended artworks with images and metadata.}
  \label{fig:system}
\end{figure}

% ─────────────────────────────────────────────────────────────────────────────
\section{Deployment Considerations}
\label{sec:deployment}

\paragraph{Latency.}
BM25 retrieval takes $<$50\,ms per query.
SBERT embedding-based retrieval takes $\sim$200\,ms (CPU) / $<$50\,ms (GPU).
GPT-5.2 re-ranking adds 0.5--2\,s of network latency.

\paragraph{Cultural coverage.}
Our bias analysis (\S\ref{sec:cultural_bias}) shows that current LLMs are
significantly less reliable for non-Western collections.
We recommend that real deployments include culturally-aware override mechanisms
and domain expert review for non-Western exhibitions.

\paragraph{Metadata quality.}
Curation performance is strongly correlated with metadata richness.
Objects with detailed \textit{culture}, \textit{medium}, and \textit{date}
fields consistently receive more accurate recommendations.
Museums should invest in structured metadata standardization (e.g., CIDOC CRM)
to maximize system utility.

\paragraph{Human-in-the-loop.}
Our informal evaluation with 3 domain experts (museum professionals and art
history graduates) confirmed that the system is most
effective as a \textit{first-pass} filter, reducing manual search time by
approximately 40\%, with final selection remaining under human control.

% ─────────────────────────────────────────────────────────────────────────────
\section{Conclusion}
\label{sec:conclusion}

We presented ExhibitionBench, the first benchmark for evaluating LLMs on
museum curation tasks.
Our experiments show that GPT-5.2 zero-shot substantially outperforms
retrieval baselines on both MEIP (Acc@1: 0.806 vs.\ BM25's 0.636) and TES
(NDCG@10: 0.678 vs.\ BM25's 0.408), establishing a strong LLM upper bound.
Yet significant performance gaps persist across cultural regions---with
GPT-5.2 0-shot scoring 0.736 Acc@1 on the dominant Western group but 0.000
on Middle Eastern objects.
We identified three systematic failure patterns (temporal mismatch: 7--18\%;
cultural drift: 19--28\%; style conflict: $\leq$1\%) that point to concrete
directions for improvement.
Our curation assistance prototype demonstrates the viability of human-in-the-loop
deployment and the utility of the benchmark for real-world museum applications.

We release ExhibitionBench, all baselines, and the demo system to encourage
the community to develop more culturally equitable museum AI.

% ─────────────────────────────────────────────────────────────────────────────
\section*{Ethics Statement}

\paragraph{Data licensing.}
Met Museum objects are licensed under CC0 (public domain).
Europeana metadata follows Creative Commons licenses (CC0 or CC BY-SA),
with attribution preserved in our dataset.

\paragraph{Cultural sensitivity.}
We recognize that algorithmic curation tools must not be used to marginalize
non-Western collections further.
Our bias analysis is explicitly designed to surface and document this risk.

\paragraph{Human evaluation.}
Any future human evaluation of this benchmark should follow IRB protocols
appropriate to the institution.

\section*{Acknowledgments}
[Blinded for review]

% ─────────────────────────────────────────────────────────────────────────────
\bibliography{references}
\bibliographystyle{acl_natbib}

\end{document}
